\section{A signed direct certificate in 519 integer variables}
\label{sec:direct519}

The principal unrestricted-integer endpoint is
\[
  P_{519}(z_1,\ldots,z_{519})=\sum_{j=1}^{321}R_j(z)^2,
\]
where every $R_j$ is the literal integral quadratic listed in the
accompanying equation file.  Its exact degree is four, its collected
expansion has $6253$ nonzero monomials, and its largest absolute
coefficient is $411060$.  There are no implicit sign conditions,
variable exponents, unspecified arithmetic ports, or program parameters
in this polynomial.  Its ordinary mathematical correctness proof gives
\[
 \mathrm{ZFC}\vdash
 \bigl[\forall z\in\mathbb Z^{519}\;P_{519}(z)\ne0\bigr]
 \longleftrightarrow\operatorname{Con}(\mathrm{ZFC}).
\]
This is an equivalence, not a proof of consistency.  No new complete
Lean proof or object-language derivation is asserted.

\subsection{The signed arithmetic interface}

Start with the square-shell reciprocal prime-power source in the preceding
sections, retaining the large carrier
\[
 M=M_0+1=((t+1)^t)^{2^{22}},\qquad
 R=2+t+x+y+M_0,\qquad K=2R.
\]
Let $F$ denote the proposed factorial, $\rho=Ft^2$,
$V=(\rho+1)^t$, $B=\binom{F^2+t}{t}$, $N=B^K$, and
$q=N(N+V+1)^K$.  Until the arithmetic graphs are justified these are
names for their proposed outputs, not assumptions about their meanings.

Replace the common Pell budget by
\[
 W=2+\sum_{a\in\mathcal A}a^2,
\]
where $\mathcal A$ is the finite set of distinct actual base, exponent,
and output atoms of the nineteen power requests.  The atom list and its
affine expansion certificate are supplied.  Thus $W\ge2$ and $|a|<W$
without a sign premise on any power output.  Retain the shared norm
\[
 A^2-W^3(W+2)(z+1)^2=1,\qquad A,z\ge0.
\]
Pell classification at base $W+1$, together with
$Y_j(W+1)\equiv j\pmod W$, gives
$A\ge(W+1)^W$.  Therefore, for $1\le n<W$,
\[
 T_n=2An-n^2-1>W^W+W>0.
\]

The finite indexed-Pell certificate inherited from the prime-power
construction admits signed comparison coordinates and signed principal
abscissa.  Its hypotheses are $A\ge2$, comparison base at least two,
$|U|\ge2$, $1\le k\le Y$, the three norm equations, the divisibility
$Y^2\mid V_{\rm aux}$, and its four congruences.  It gives
\[
 |X|=X_k(A),\qquad Y=Y_k(A).
\]
For completeness, the signed extension follows by classifying absolute
values of all three norm solutions.  The abscissa congruence gives the
comparison index congruent to either sign of the principal index modulo
$2n$, and the ordinate congruence gives either sign of $k$ modulo $4Y$.
Squared divisibility gives $Y\mid n$.  Consequently the principal index
is congruent to $k$ or $-k$ modulo $2Y$.  Both lie in $[1,Y]$, so they
coincide, including the endpoint case.  This is the same indexed-Pell
lemma used by the inherited 553-integer proof, not a finite-test inference.

For the signed power row
\[
 X+(n-A)Y=m+T_n u
\]
with $1\le n,k<W$ and $|m|<W$, a positive $X$ yields
$m\equiv n^k\pmod {T_n}$ and hence $m=n^k$ by the strict interval bound.
A negative $X$ instead gives $mn^k\equiv-1\pmod {T_n}$: multiply the two
conjugate expressions and use the norm.  Since
$|mn^k+1|<T_n$, this forces $mn^k=-1$.  It is impossible for $n\ge2$;
for $n=1$ it allows only $m=-1$.  Thus a single request with base at least
two orients the shared index pair for all requests at that exponent.

\subsection{Why the auxiliary nonzero guard disappears}

The earlier 553-integer construction inserted a three-witness guard for
$U^2\ge2$.  It is unnecessary in the present reciprocal source.
Every seed exponent is already positive before any power has been
interpreted.  The complete list consists of
\[
 t,\ K,\ Ft^2+t,\ F^2+t,\ y+1,\ x+1,
 \quad\hbox{and the positive encoded representatives.}
\]
The seed $t+1$ is derived by the existing addition identity and is positive
as well.  The retained seed guards imply $t\ge1$, $x,F,r_0\ge0$ and
positive encoded representatives.  The final carrier squaring already
implies $M_0+1\ge0$, whether or not its power input has the intended value.
Hence
\[
 y=M_0+r_0+2\ge1,\qquad
 R=t+x+r_0+2(M_0+1)+2\ge3.
\]
Every seed inequality $Y_i=k_i+e_i$, $e_i\ge0$, therefore gives $Y_i\ge1$.
The unchanged common auxiliary product
\[
 V_{\rm aux}=\prod_iY_i^2\,(v+1),\qquad v\ge0,
\]
is positive.  Its norm yields $U^2\ge A^2\ge4$, proving the omitted
guard without a circular use of any power output.  The emitted certificate
expands each seed exponent after substituting
$M_0=Z-1$, $y=Z+r_0+1$ with $Z\ge0$; every resulting polynomial has
nonnegative coefficients and positive constant term.

\subsection{The noncircular recovery order}

Sixty natural seed coordinates are guarded by
$s=a^2+b^2+c(c+1)$.  Legendre's three-square theorem applied to $4s+1$
proves that this represents exactly the nonnegative integers.  The
coordinate list and an ordered exact sign-inference certificate are part
of the artifact; all other coordinates are unrestricted.

Let $D$ be the factorial binomial output, and let $r,s,o$ be its retained
nonnegative remainders and slacks.  Before interpreting any binomial,
the actual raw rows imply
\[
 D=r+s+1\ge1,\qquad q=D+o+1\ge2.
\]
The existing request $q^t$ therefore orients the $t$-indexed pair.  The
request $(t+1)^t$ now has its true value, recovering $M_0\ge0$ and the
usual endpoint signs.  The request $(Ft^2+1)^t$ is at a positive base,
possibly initially one, and shares the already oriented pair; hence
$V\ge1$.  The raw factorial rows give $DF=V+s$, so $F\ge1$ and then
$V\ge2$.

The reciprocal index row is
\[
 C+F^2+1=B i_q.
\]
Here $C\ge1$ and the output $B$ is retained nonnegative, so $B>0$.
The request $B^K$ either has its intended positive output $N$, or its
only remaining inverse case has $B=1$ and $N=-1$.  In both cases the
second $K$-indexed base satisfies
\[
 N+V+1\ge2.
\]
That request orients the $K$ pair and excludes the inverse case.  Thus
$N=B^K>0$ and the radix has its intended value.  Every remaining seed
has an actual request at base $q$ or $q+1$, both at least two.  All its
exponents are positive by the preceding subsection.  This recovers all
nineteen powers, without first assuming the factorial, product-binomial,
or range-binomial conclusions under validation.

The natural source's ordered factorial, low-digit, prime-power lifting,
and CRT arguments now apply unchanged.  Required Euclidean quotients
are nonnegative automatically from their positive divisors and bounded
nonnegative remainders.  The auxiliary congruence quotients and comparison
coordinates need not become nonnegative; only the proved Pell congruences
use them.  Simultaneous Pell completeness supplies compatible witnesses
for any genuine semantic certificate.

\subsection{Factoring the Horn squares over the integers}

Write the nine ordinary axiom balances, including truth, and the two
inference balances as $L_1,\ldots,L_{11}$.  Let $D_b$ be the guarded
bundle balance and
\[
 G=(c-a)^2(c-b)^2-(d+1)^2.
\]
The original Horn summand is
\[
 \left(\prod_{i=1}^{11}L_i^2\right)(D_b^2+G^2).
\]
Instead form signed intermediate products and one final square:
\[
 H=\prod_{i=1}^{11}L_i,
 \qquad H^2(D_b^2+G^2).
\]
These are identically equal over every commutative ring.  Likewise form
$h=(c-a)(c-b)$ and $G=h^2-(d+1)^2$ before squaring $G$.
This removes duplicated square gates.  The two beta squares remain
outside the Horn product.  Merely replacing $H^2$ by $H$ in their common
sum would be unsound because of cancellation; that change is not made.

The factorization certificate identifies every original square root,
the eleven product links, and every changed row.  It checks that no
intermediate factor escapes this subnetwork.  Given an old zero, define
the new signed factors and products; given a new zero, restore the old
squared factors and their products.  The visible Horn value and all
other coordinates agree exactly.  Existing sign guards mention none
of the changed internal coordinates.  Therefore the change preserves
integer-zero existence in both directions, independently of arithmetic
or logical assumptions.

\subsection{Projection and integral translation}

After the sixty sign guards, the signed source has 734 coordinates and
536 rows.  Factorization has 735 coordinates and 535 rows before pruning.
The recorded monic projection deletes 214 coordinates, of which 168 have
affine and 46 have quadratic integral restorations, and removes two
unconstrained coordinates.  It leaves 519 coordinates and 321 rows.
Every deleted defining row is exactly $z-Q$ up to sign.  A quadratic
$Q$ is substituted only where $z$ occurs linearly in every other row,
so all residuals remain quadratic.  Reversing the substitutions gives
the unique integral lift of every deleted defined coordinate.

Finally translate each retained encoded seed, $t_0$, $M_0$, and retained
coordinate whose name ends in \texttt{minus1} by minus one: the convention
is $z_{\rm old}=z_{\rm new}-1$.  The exact coordinate indices, not the
names, define the transformation in its certificate.  Its inverse adds
one, so this is a bijection of $\mathbb Z^{519}$.  It absorbs offsets
that otherwise inflate constant and linear coefficients.  The degree
and variable count stay unchanged, while monomials fall from 6816 to
6253 and coefficient height from 6474801 to 411060.  This is a coordinate
choice, not a claim of globally minimal height.

\subsection{Finite evidence and logical scope}

The independent-friendly artifact consists of the raw quadratic source,
square-budget and sign certificates, source-factorization certificate,
monic elimination certificate, integral translation certificate, final
quadratic residual list, and fully collected quartic coefficients.
The standard-library reproduction script imports no downloaded builder.
Exact coefficient reconstruction checks these finite algebraic claims;
the unbounded Pell, CRT, and proof-flow arguments establish their meaning.

The numerical-calculus calibration remains
\texttt{KEndpoint.c222\_iff\_zf\_refutation} followed by
\texttt{ConstructibleChoice.zf\_refutation\_iff\_zfc\_refutation}.
Square-shell recursive recoding fixes truth and falsity and transports
all formula and variable positions.  No K215 execution, Jones polynomial,
or Sun variable reduction occurs in this construction.  Full transitive
formal-kernel checking, global minimality, and worldwide priority are
not claimed.

\subsection{The separate 517 variable tradeoff}

Applying the same signed compiler to the derived-label factorial-carrier
source yields a second quartic in 517 unrestricted integers, with 310
residuals, 6774 collected monomials, and coefficient height 3904267.
The complete arithmetic proof for that source is separate: it replaces
the large carrier by $F=t!$ and certifies the derived label sum by a
range binomial rather than a ninth independent local coordinate.
It is not a further projection of $P_{519}$.

The signed proof adapts without an additional guard.  Here
$y=F+r_0+1\ge1$ and $K=2(2+t+x+y+F)>0$ immediately.  The additional
seed exponent is the derived label sum $J$; its constructor circuit
uses only nonnegative addition and multiplication and has an explicit
positive constant summand.  Therefore $J\ge1$ before power recovery.
The emitted version-two positivity certificate follows these gates
compositionally, retaining positive lower bounds capped at 64; it
does not expand the enormous polynomial $J$.  All other seed indices
are the same positive types already listed.  Thus the product again
forces $V_{\rm aux}>0$ and $|U|\ge2$ without a guard.  The raw $q\ge2$
argument orients $t$ directly, the factorial ceiling forces $F\ge1$,
and the same two-request argument orients $K$.  Source factorization,
integral projection, and the explicitly recorded coordinate translation
then establish the 517-variable endpoint.  It improves the variable
count by two but has more terms and larger coefficients than the
primary 519-variable presentation; neither dominates the other.
